% Lösungen Einheit 4 von 4 – Integralfunktion (zusammenbau v0.1)
\einheitenkopf{Einheit 4 von 4 · Integralfunktion}

\erg{15}{a) ja – die Flächenstücke heben sich auf; das Dreieck unter der Achse ist so groß wie das darüber. \quad b) $x = 2$ (gleiche Grenzen); $J'(x) = x^2 - 1$ \quad c) $x = 1$ (gleiche Grenzen) und $x = 5$: wegen der Punktsymmetrie zu $(3 | 0)$ ist das Flächenstück von $1$ bis $3$ über der Achse so groß wie das von $3$ bis $5$ darunter. \quad d) höchstens $3$: $J$ ist ganzrational vom Grad $3$, weil $f$ den Grad $2$ hat. \quad e) Wendestelle $x = 2$, weil $f'(x) = 2x - 4$ dort das Vorzeichen wechselt; $J(2) = \frac{8}{3} - 8 = -\frac{16}{3}$ \quad f) wahr – $J(1) = 0$; links von $1$ ist $J$ negativ, weil $f$ dort positiv ist; rechts von $1$ steigt $J$ bis $3$ und fällt dann unbeschränkt, also gibt es genau eine weitere Nullstelle rechts von $3$ – nicht die Nullstelle $3$ von $f$.}
\erg{16}{a) um $9$ nach unten: $M(x) = L(x) - L(3)$, und $L(3) = 18 - 9 = 9 > 0$, weil $f$ zwischen $0$ und $3$ positiv ist.}
\erg{17}{a) Bei $x = 2$ waagerechte Tangente, weil $J'(2) = f(2) = 0$; Wendepunkt, weil $f = J'$ bei $2$ einen Tiefpunkt hat – ein Terrassenpunkt, kein Extrempunkt, da $f$ dort das Vorzeichen nicht wechselt.}
\erg{18}{a) Ole hat die Nullstelle des Integranden genommen; richtig: $J(x) = \frac{1}{2}x^2 - 2x = 0$ für $x = 0$ und $x = 4$. \quad b) Bei gleichen Grenzen ist das Integral $F(4) - F(4) = 0$ für jede Stammfunktion $F$ von $f$.}
